\subsection{Critical graph budgets can coexist with an easy projection}
\label{sec:kernel-realization}

The critical gate and cycle budgets alone do not force a hard
projection. A large class of cubic graphs can occur as exact residual
kernels with nearly critical budgets and a simple projection. The
distinction between the supplied circuit and a minimum circuit is
essential here.

\begin{theorem}[Cubic realization with controlled budgets]
\label{thm:kernel-realization}
Let $K$ be a finite simple 2-vertex-connected cubic graph and put
$k=|V(K)|/2+1$. Set $t=\lfloor k/3\rfloor+1$; when $K$ is
triangle-free, one may instead take $t=\lceil k/3\rceil$.
There is an acyclic AND-only verifier on $t$ ordinary and $t$ witness
inputs with
\[
 (p,q,h,b,\ell,\kappa)
   =(0,k+2t-1,t,t,2k+3t-2,k).
\]
It has no constants, unary gates, repeated input wires at a gate,
unused inputs, or gates outside its output cone. A specified legal
equality splitting and reduction sequence produces exactly $K$, with
zero cycle-rank loss. Every input is essential, but
\[
 f(x,y)=\AND_{i=1}^{t}x_i\land\AND_{j=1}^{t}y_j,
 \qquad g(x)=\AND_{i=1}^{t}x_i,
 \qquad \C(f)=2t-1,\quad \C(g)=t-1.
\]
Thus $h+b+\kappa=q+1$ and $t=k/3+O(1)$, while the supplied
circuit exceeds its minimum size by exactly $k$ gates.
\end{theorem}

The syntactic conditions in the statement are those used by the
compiler's normalization. They do not prohibit identities involving
several gates, such as associativity and idempotence of AND.
We give the construction and its accounting in full.

\begin{proof}
First give $K$ an acyclic orientation with one source and one sink.
This is the classical $st$-numbering property of 2-connected graphs
\citep{EvenTarjan1976}. An existence proof suffices here. Start with
a cycle, choose adjacent vertices $s,z$, and order the vertices along
the $s$-to-$z$ path obtained by deleting the edge $sz$. If some
vertices remain outside the current set, one of their components has
two distinct attachments to the current set; otherwise its unique
attachment would be a cut vertex. Choose a path through that component
between two such attachments $u,v$, with internal vertices outside.
Orient the path from its earlier to its later endpoint and insert its
internal vertices, in order, immediately after the earlier endpoint.
Every old internal vertex retains neighbors on both sides, and each
new vertex has its two path neighbors on opposite sides. Repeat until
all vertices are ordered and orient every edge forward. Only $s$ is
a source and only $z$ is a sink.

\emph{The circuit skeleton.}
Every other vertex is a fork with indegree one and outdegree two,
or a join with indegree two and outdegree one. If their numbers are
$F,J$, the degree sum gives
\[
 3+F=J+3,\qquad F+J+2=2(k-1),\qquad F=J=k-2.
\]
Put witness $y_1$ at $s$. A fork copies its incoming wire at no gate
cost. A join ANDs its incoming wires. At $z$, two AND gates combine
the three incoming wires. This skeleton has $k$ gates.
Temporarily assign each join output its own formal wire name, even
if that join receives the same name twice. Call such a join a conflict.

The sink receives at least two different names. Indeed a topologically
last join exists, and any path from it to the sink contains only forks.
If its output passes through a nonempty fork tree, a last fork would
have both outgoing edges to the sink, contradicting simplicity.
Its output therefore reaches the sink directly and on just one edge.
It cannot name all three sink inputs. Choose two distinct names for
the first sink gate; its new output and the third original input
are distinct at the second gate.

\emph{The conflict budget.}
The subgraph on the source and forks is a forest. A cycle with an
acyclic orientation would have a vertex with two incoming edges,
which these vertices cannot have. Its components are exactly the
trees carrying one formal wire name. A component not containing $s$
with $u$ forks has one incoming edge and $u+1$ outgoing edges. Each
conflict it supplies uses two outgoing edges. For $u=1$ there is no
conflict, by simplicity; for $u\ge2$ their number is at most
\[
 \left\lfloor\frac{u+1}{2}\right\rfloor\le\frac{2u}{3}.
\]
The inequality is checked directly at $u=2$ and follows from
$(u+1)/2\le2u/3$ for $u\ge3$.

The source component with $u$ additional forks has $u+3$ outgoing
edges. At $u=0$ it has no conflict; at $u=1,2$ it has at most two.
For $u\ge3$ its number of conflicts is at most
$\lfloor(u+3)/2\rfloor\le2u/3+4/3$.
The latter bound covers the small cases as well. Summing over all
$F=k-2$ forks bounds the total number $R$ of conflicts by
\[
 R\le\lfloor2k/3\rfloor.
\]
If $K$ is triangle-free, the $u=1$ source component has no conflict,
since one would make a triangle through the source, fork, and target
join. The source bound improves to $2u/3+1$, giving
\[
 R\le\lfloor(2k-1)/3\rfloor
 \qquad\text{in the triangle-free case}.
\]

There are $2t-1$ unused inputs available: the $t$ ordinary inputs and
$t-1$ additional witnesses. The chosen $t$ ensures $R\le2t-1$ in
each case. On one incoming edge of each conflicting join insert
\[
 \text{new wire}=\text{old wire}\land\text{fresh input},
\]
using a different input each time. Insert any remaining inputs by
such gates in a chain on one original edge. The fresh output names
split copied wires; they never identify distinct old names. Hence
all conflicts are repaired without creating new ones, including at
the two sink gates. The orientation gives a topological gate order.

\emph{The verifier budget and its function.}
Each attachment has a witness-dependent old input, so all gates are
witness-dependent. There are $k+2t-1$ gates, every ordinary input
appears at one gate pin, and all $t$ witnesses are used. Therefore
$p=0$, $h=b=t$, and
\[
 \ell=2q-t=2k+3t-2,\qquad
 \kappa=\ell-q-b+1=k.
\]
Every vertex reaches the unique sink. ANDing along the circuit thus
gives the conjunction of all inputs, proving the displayed $f,g$.
Each input is essential by setting all other inputs to one. The
effective-input lower bound in Lemma~\ref{lem:inputs},
met by an AND chain, gives $\C(f)=2t-1$ and $\C(g)=t-1$.

\emph{The exact residual kernel.}
Form the actual consistency incidence graph. A wire with $u$ forks
after its defining join has $u+1$ uses and one defining incidence.
Choose its equality splitting to be precisely that fork tree: an
equality of degree $u+2$ is replaced by $u$ ternary equalities.
For $u=0$, keep its degree-two equality. The source wire with $u$
forks has $u+3$ uses; choose its splitting into the source and those
$u$ forks. Attachments divide these trees into smaller pieces, to
which the same construction applies. Reversing these splittings
contracts each equality tree to the actual incidences of one circuit
wire, verifying that they are permitted choices in
Lemma~\ref{lem:cubic-reduction}.

Contract the output-one assertion through the output equality into
the second sink gate. That constraint now has two ports. Contract
the degree-two equality between the sink gates and then the remaining
two-port constraint into the first sink gate. Its surviving table
is the conjunction of the three original incoming bits, at vertex $z$.
For each additional witness, contract its degree-one equality into
its attachment constraint, leaving two ports. For example,
$\OR_y[v=u\land y]=[v\le u]$: the resulting table is retained,
not replaced by an equality. Ordinary attachments
already have two ports, since their ordinary inputs are coefficients.
Contract these two-port constraints and all degree-two wire equalities
along their original edges. Their updated neighboring tables retain
the semantic effects of the attachments.

Topologically these operations undo edge subdivisions, the sink tree,
and attached leaf trees. They leave exactly the marked vertices and
edges of $K$. No loop or parallel-edge reduction is needed: every
step is a permitted single-edge contraction and preserves cycle rank.
Since $K$ is simple and cubic, the reduction terminates there, with
the stated zero loss and the linear operation cost of the compiler.
\end{proof}

For triangle-free $K$ with $k=3t$, the parameters are exactly
\[
 (p,q,h,b,\ell,\kappa)=(0,5t-1,t,t,9t-2,3t).
\]
Thus acyclicity, essential inputs, the critical budget, and zero cycle
loss do not prevent these kernel topologies. The theorem asserts
existence of the specified reduction, not that every reduction must
return $K$. It does not realize the hardness hypothesis of
Theorem~\ref{thm:critical-structure}: its supplied circuit is
nonminimal and its projection is easy. Stronger semantic reductions
or choosing a smaller equivalent verifier can remove the apparent
structural difficulty.
